\section{The unmasked Rayleigh rate and nonnegative controls}
\label{S:unmasked-variational}

Fix a nonzero covariance satisfying Hypothesis~\ref{H:homog}, including its
original amplitude, and set $a = \alpha_0$. Define the full class
\begin{equation}\label{E:variational-positive-class}
  \mathcal{C}_{\mathrm{pos}}
  =\{\Psi\in C_c^\infty((0,1)\times\R^d;\R):
  \Psi\ge0,\ \Psi\ne0\}.
\end{equation}

\begin{figure}
  \centering
  \begin{tikzpicture}[x=1cm,y=1cm,font=\normalsize]
    % Identical panel widths, baselines, box sizes and type sizes.
    \foreach \shift in {0,7.6} {
    \begin{scope}[xshift=\shift cm]
      \draw[ham/arrow] (.85,0)--(6.4,0) node[right] {$x$};
      \draw[ham/arrow] (.85,0)--(.85,2.3) node[above] {$s$};
      \node[anchor=east] at (.7,.6) {early};
      \node[anchor=east] at (.7,1.75) {late};
      \node at (2.1,-.3) {left};
      \node at (4.9,-.3) {right};
      \filldraw[draw=hamblue,fill=hamblue!12]
      (1,.2) rectangle (3.2,1);
      \filldraw[draw=hamteal,fill=hamteal!12]
      (3.8,1.35) rectangle (6,2.15);
      \node at (2.1,.6) {$>0$};
      \node at (4.9,1.75) {$>0$};
    \end{scope}
    }
    \node[font=\normalsize\bfseries] at (3.5,3.5) {Single product};
    \node[font=\normalsize\bfseries] at (11.1,3.5) {Two-bump sum};
    \node at (3.5,2.95) {$\Psi=\phi\psi$};
    \node at (11.1,2.95) {$\Psi=\phi_1\psi_1+\phi_2\psi_2$};
    \foreach \x/\y in {4.9/.6,2.1/1.75} {
    \filldraw[draw=hamgold,dashed,fill=hamgold!12]
    (\x-1.1,\y-.4) rectangle (\x+1.1,\y+.4);
    \node[text=hamgold] at (\x,\y) {$>0$ forced};
    }
    \node at (12.5,.6) {$0$};
    \node at (9.7,1.75) {$0$};
    \node[align=center,text width=6.6cm] at (3.5,-.95)
    {Four positive blocks:\\ the cross blocks are forced.};
    \node[align=center,text width=6.6cm] at (11.1,-.95)
    {Two positive blocks:\\ the cross blocks can be zero.};
    \node at (7.3,-1.65)
    {Schematic supports, not an optimizer or a density claim.};
  \end{tikzpicture}
  \caption[Controls that need not factor]{The full control class allows spatial
    profiles to change with time, as in this sum of two smooth nonnegative product bumps that no single
    time--space product can represent. The full class is infinite-dimensional;
    see \ProofRef{Theorem}{T:intrinsic-second-moment}.}
  \label{Fig:nonseparable-controls}
\end{figure}

It includes nonseparable tests. It is nonempty: if $b(z) = e^{-1/(1-z^2)}$ on
$|z| < 1$ and zero otherwise, then $b(4s-2)\prod_{j=1}^db(x_j)$ belongs to this
class. All boundary derivatives vanish because exponential decay dominates every
power of $(1-z^2)^{-1}$. For each $\Psi$ retain the original mean $h_{\Psi,s}$
and original variance $v_\Psi = \E W(\Psi)^2$ from~\eqref{E:compact-raw-mean}--%
\eqref{E:compact-raw-cost}. Neither the test nor its covariance is changed by a
physical mask. We write $h_s$ when $\Psi$ is fixed.

\subsection{Global bounds and the unmasked spectral bottom}

We first record a bound for the covariance smoothed by the heat kernel. The
proof of Lemma~\ref{L:variational-global-mean} uses it to bound
$\Gamma*|D\Psi(t,\cdot)|$ uniformly in $t,x$.

\LeanMarker{L:variational-heat-envelope}
\begin{lemma}[Gaussian envelope for the original covariance]
  \label{L:variational-heat-envelope}
  For the normalized heat kernel $p_\epsilon(x) =
  (2\pi\epsilon)^{-d/2}e^{-|x|^2/(2\epsilon)}$,
  \begin{equation}\label{E:variational-heat-envelope}
    (\Gamma*p_\epsilon)(x)
    =\int e^{i\xi\cdot x}e^{-\epsilon|\xi|^2/2}\mu(\ud\xi),
    \qquad
    0\le(\Gamma*p_\epsilon)(x)
    \le(\Gamma*p_\epsilon)(0)<\infty.
  \end{equation}
  The value at zero is strictly positive.
\end{lemma}

\begin{proof}
  We first use temperedness to make all Gaussian pairings finite. We identify
  the convolution and spectral expressions as continuous distributions, then
  use positivity to obtain the bound by the value at zero.


  \medskip
  \noindent\textit{Step 1.~}We obtain integrable majorants for the Gaussian
  weights.
  A positive tempered Radon measure has polynomial ball growth. Indeed, evaluate
  its tempered bound on a nonnegative smooth cutoff equal to one on $B_1$,
  dilated by $R \ge 1$. Its finitely many Schwartz seminorms are bounded by
  $CR^N$, so positivity gives a ball-mass bound $CR^N$. Summing dyadic annuli
  shows that $(1+|x|)^{-L}$ is integrable for $L > N$. Apply this separately to
  $\Gamma$ and $\mu$. Gaussian weights are therefore integrable against both
  measures. The classical convolution on the left is finite and smooth locally
  in $x$: on compact $x$ sets its translated Gaussian derivatives are bounded by
  any fixed sufficiently fast polynomial weight in the integration variable.

  On the annulus $2^j\le |x|<2^{j+1}$, the mass bound contributes at most a
  constant times $2^{jN}$ and the decreasing weight contributes at most a
  constant times $2^{-jL}$. Their product is summable when $L>N$. A Gaussian
  and each of its derivatives decrease faster than any such polynomial weight.

  \medskip
  \noindent\textit{Step 2.~}We identify the two continuous expressions and
  bound their common value.
  Pair the convolution with a Schwartz test $\varphi$. Absolute Fubini is
  justified by these weights. The resulting pairing is that of $\Gamma$ with
  $p_\epsilon*\varphi$. Applying the spectral convention in
  Hypothesis~\ref{H:homog} and the identity $\int p_\epsilon(y)e^{-i\xi\cdot
  y}\ud y = e^{-\epsilon|\xi|^2/2}$ gives exactly the right side
  of~\eqref{E:variational-heat-envelope} as a distribution. The heat kernel is
  even, so either Fourier sign gives that same Gaussian multiplier. Its spectral
  integral is absolutely finite and continuous in $x$ by dominated convergence.
  Thus, equality of the two continuous distributions gives equality at every
  $x$, without an extra normalization factor. The physical convolution is
  nonnegative. Evenness of $\mu$ makes the spectral integral real, and its
  absolute value is bounded by its total Gaussian mass, the value at zero.
  Finally, $\Gamma \ne 0$ and $p_\epsilon > 0$ give strict positivity there.
  The spectral phase has modulus one, so the absolute value of its integral is
  at most the integral of the positive Gaussian weight. At $x=0$ that phase is
  one. Physical positivity makes the common value nonnegative, completing the
  two-sided comparison.
  This proves Lemma~\ref{L:variational-heat-envelope}.
\end{proof}

\LeanMarker{L:variational-global-mean}
\begin{lemma}[Global regularity of the original mean]
  \label{L:variational-global-mean}
  Every mixed time and space derivative of $h_s$ is globally bounded in space,
  uniformly on compact time intervals. In particular, for $I = [0,1]$ there are
  finite constants $M,C_\Psi$ such that
  \begin{equation}\label{E:variational-global-mean}
    0\le h_s(x)\le M,\qquad
    \|h_s-h_t\|_\infty\le C_\Psi|s-t|,
    \qquad s,t\in I.
  \end{equation}
\end{lemma}

\begin{proof}
  Extend $\Psi$ by zero in time; its compact support inside $(0,1)$ makes this a
  smooth extension. For any mixed derivative $D\Psi$, a fixed $\epsilon > 0$ and
  compactness of its common spatial support give $|D\Psi(t,x)| \le
  C_Dp_\epsilon(x)$ for all $t,x$. Here the heat kernel has a positive minimum
  on that support. Lemma~\ref{L:variational-heat-envelope} consequently bounds
  $\Gamma*|D\Psi(t,\cdot)|$ uniformly in $t,x$. Ordinary local differentiation
  against the Radon measure gives the corresponding derivatives of the
  convolution. After the change $r = s-t$,
  \[
    h_s(x)=\int_\R |r|^{-a}
    (\Gamma*\Psi(s-r,\cdot))(x)\ud r.
  \]
  The extension of the test is supported in one bounded time interval. For $s$
  in a fixed compact interval, the condition that $\Psi(s-r,\cdot)$ is nonzero
  therefore confines $r$ to another fixed bounded interval. This confinement
  also holds for every derivative of the test.
  For $s$ in a compact interval the relevant $r$ lie in a fixed bounded
  interval. The preceding global bound times $|r|^{-a}$ is integrable since $a <
  1$. Every derivative thus falls on the smooth test and passes through the
  integral by dominated convergence, with the claimed global spatial bounds.
  Nonnegativity follows from positive kernels. The bounded time derivative
  proves the Lipschitz estimate. No directional decay of the mean has been
  inferred. This proves Lemma~\ref{L:variational-global-mean}.
\end{proof}

Lemma~\ref{L:variational-global-mean} shows that $h_s$ is bounded by $M$ for $s
\in I$ and is Lipschitz in $s$ in the supremum norm. We now use these two bounds
to define the rate $\lambda_{\Psi,s}$ of the potential $h_s$ without a mask and
to show that it is Lipschitz on $I$.

\LeanMarker{P:variational-unmasked-rate}
\begin{proposition}[The unmasked Rayleigh rate]
  \label{P:variational-unmasked-rate}
  On complex $L^2(\R^d)$, the operator $H_s = -\Delta-h_s$ is self-adjoint on
  $H^2$, with form domain $H^1$ and form $q_s(f) = \|\nabla f\|_2^2-\int
  h_s|f|^2$. Its spectral bottom $E_s$ is uniformly negative on $I$. The
  positive rate $\lambda_{\Psi,s} = \sqrt{-E_s}$ satisfies
  \begin{equation}\label{E:variational-real-rayleigh}
    \lambda_{\Psi,s}^2
    =\sup_{\substack{f\in H^1(\R^d;\R)\\\|f\|_2=1}}
    \left(\int h_sf^2-\|\nabla f\|_2^2\right).
  \end{equation}
  It is Lipschitz on $I$, and $I_\Psi: = \int_0^1\lambda_{\Psi,s}\ud s$ is
  positive and finite. No eigenstate at the unmasked bottom is asserted.
\end{proposition}

\begin{proof}
  We construct the unmasked self-adjoint operator using bounded perturbation
  and identify its form domain. The spectral theorem gives the real Rayleigh
  formula; the common compact trial makes its bottom uniformly negative and
  its square-root rate Lipschitz.


  \medskip
  \noindent\textit{Step 1.~}We prove self-adjointness on the stated operator
  domain.
  The Fourier Laplacian $A = -\Delta$ is self-adjoint on $H^2$, being
  multiplication by $|\xi|^2$ with its maximal domain. Multiplication $B_s =
  -h_s$ is bounded self-adjoint with norm at most $M$. Thus, $H_s = A+B_s$ is
  densely defined, symmetric and closed on $H^2$; the two graph norms are
  equivalent under a bounded perturbation. For real $b > M$,
  \[
    H_s-ib=(I+B_s(A-ib)^{-1})(A-ib).
  \]
  The first factor has a Neumann inverse because $\|B_s(A-ib)^{-1}\| \le M/b <
  1$, and the same holds with $+ib$. To check self-adjointness directly, for $u
  \in D(H_s^*)$ solve $(H_s-ib)v = (H_s^*-ib)u$ with $v \in D(H_s)$. Then $u-v
  \in \ker(H_s^*-ib) = \operatorname{Ran}(H_s+ib)^\perp = \{0\}$. Hence,
  $D(H_s^*) = D(H_s)$.

  Surjectivity at both imaginary spectral parameters is the key point in the
  adjoint-domain argument. Solving at $-ib$ produces $v$, while surjectivity
  at $+ib$ makes the orthogonal complement of its range zero. Thus, the vector
  $u-v$ must vanish, giving the reverse domain inclusion.

  \medskip
  \noindent\textit{Step 2.~}We identify the form domain by equivalent
  completions.
  Set $c_{\mathrm{shift}} = 1+M$ and $J_s = H_s+c_{\mathrm{shift}} \ge 1$.
  Integration by parts on $H^2$ gives
  \[
    \|J_s^{1/2}f\|_2^2=q_s(f)+c_{\mathrm{shift}}\|f\|_2^2,\qquad
    \|f\|_{H^1}^2\le\|J_s^{1/2}f\|_2^2
    \le c_{\mathrm{shift}}\|f\|_{H^1}^2.
  \]
  Spectral cutoffs $1_{[1,n]}(J_s)$ show that $H^2 = D(J_s)$ is a core for the
  half-power norm; Fourier cutoffs give its density in $H^1$. Completion in the
  equivalent norms identifies $D(J_s^{1/2})$ and $H^1$ as $L^2$ subspaces, since
  both norms dominate $L^2$. This proves the stated form domain.


  \medskip
  \noindent\textit{Step 3.~}We obtain the real Rayleigh formula without an
  eigenstate assumption.
  The spectral lower bound gives one direction of the Rayleigh identity. For the
  reverse, nonzero spectral projections in $[E_s,E_s+\epsilon]$ give unit $H^2$
  vectors with Rayleigh values at most $E_s+\epsilon$. The bottom is finite
  because the operator is bounded below and has finite Rayleigh tests. Let
  $\epsilon\downarrow0$. Real tests give the same infimum: $q_s(u+iv) =
  q_s(u)+q_s(v)$, so the quotient is the norm-weighted average of the nonzero
  real quotients. This proves the Rayleigh identity with $-E_s$ on the left. The
  following negative upper bound justifies writing that value as
  $\lambda_{\Psi,s}^2$ in~\eqref{E:variational-real-rayleigh}.

  A vector with spectral support in the bounded interval $[E_s,E_s+\epsilon]$
  belongs to the operator domain, because the spectral multiplier is bounded
  on that support. Its normalized spectral measure has total mass one, so its
  Rayleigh quotient lies in the same interval. For a complex vector, the
  numerator and squared norm split into real and imaginary parts; at least one
  nonzero real part has quotient no greater than their weighted average.

  \medskip
  \noindent\textit{Step 4.~}We obtain strict negativity and uniform rate
  continuity.
  Proposition~\ref{P:compact-original-common-trial} supplies, for this original
  $\Psi$, one unit compactly supported $w$ and a finite plateau mask with a
  common strictly negative upper bound $\epsilon_0 > 0$. Its trial has support
  in $B_{2L}$; choose its plateau radius $R_0 \ge 2L$. The unmasked form on this
  same $w$ is no larger. Consequently, $-M \le E_s \le -\epsilon_0 < 0$, and
  $\sqrt{\epsilon_0} \le \lambda_{\Psi,s} \le \sqrt M$. Comparing identical unit
  form tests gives $|E_s-E_t| \le \|h_s-h_t\|_\infty$ and therefore
  \[
    |\lambda_{\Psi,s}-\lambda_{\Psi,t}|
    \le \frac{C_\Psi}{2\sqrt{\epsilon_0}}|s-t|.
  \]
  This also proves positivity and finiteness of $I_\Psi$. No compact-potential
  spectral theorem was applied to $H_s$. This proves
  Proposition~\ref{P:variational-unmasked-rate}.
\end{proof}

\subsection{Removing masks from the numerical lower rate}

Proposition~\ref{P:raw-block-original-lower} gives the lower bound with the rate
of the masked potential $m_Rh_s$, for one fixed mask. We now let the mask radius
$R$ tend to infinity and replace this rate by the rate $\lambda_{\Psi,s}$ of
Proposition~\ref{P:variational-unmasked-rate}.

\LeanMarker{P:variational-remove-mask}
\begin{proposition}[Monotone rates and the original point moment]
  \label{P:variational-remove-mask}
  There is an increasing family of admissible radial masks $m_R$ tending to one,
  for which the compact ground rates satisfy
  $\lambda_{R,s}\uparrow\lambda_{\Psi,s}$ uniformly on $I$ as $R \to \infty$,
  with the compact ground statements used for $R \ge R_0(\Psi)$. For the
  original canonical solution and any fixed $x$,
  \begin{equation}\label{E:variational-unmasked-point-lower}
    \liminf_{T\to\infty}T^{-\kappa}\log\E|u(T,x)|^2
    \ge2I_\Psi-\frac{v_\Psi}{2}.
  \end{equation}
\end{proposition}

\begin{proof}
  We choose nested physical masks, compare their Rayleigh suprema with the
  unmasked supremum, and prove uniform convergence in the time parameter. We
  then take a numerical supremum of already established fixed-mask bounds.


  \medskip
  \noindent\textit{Step 1.~}We construct the nested masks.
  Put $F(t) = e^{-1/t}$ for $t > 0$ and zero otherwise, and define
  \[
    \chi(r)=\frac{F(2-r)}{F(2-r)+F(r-1)},\qquad
    m_R(x)=\chi(|x|/R).
  \]
  The denominator is positive. The numerator decreases and the second
  denominator term increases; thus $\chi$ is nonincreasing, equals one on $r \le
  1$, and zero on $r \ge 2$. Flatness of $F$ at zero gives smoothness. Since
  $\chi$ is constant near zero, $m_R$ is smooth also at the spatial origin.
  These even $[0,1]$ masks increase pointwise to one.


  \medskip
  \noindent\textit{Step 2.~}We pass the Rayleigh values and rates to their
  unmasked limits.
  Choose $R_0$ so the preceding common trial lies in the plateau. Every $R \ge
  R_0$ has the same negative margin. For each fixed $R$, the compact-family
  theorem applies to $m_Rh_s$. Its real Rayleigh formula follows from the same
  real/complex quotient argument as~\eqref{E:variational-real-rayleigh}. The
  rates are continuous, positive, at most $\sqrt M$ and monotone in $R$. For
  every fixed real unit $f \in H^1$, dominated convergence gives $\int m_Rh_sf^2
  \to \int h_sf^2$. Hence,
  \[
    \sup_{R\ge R_0}\lambda_{R,s}^2
    =\sup_{\|f\|_2=1,\ f\in H^1(\R^d;\R)}
    \sup_{R\ge R_0}
    \left(\int m_Rh_sf^2-\|\nabla f\|_2^2\right)
    =\lambda_{\Psi,s}^2.
  \]
  For each fixed test, $0\le m_Rh_sf^2\le M f^2$, an integrable majorant. Its
  potential integral therefore converges to the unmasked one. Taking the
  supremum over the product of the radius and test index sets in either order
  gives the displayed equality; no maximizer is required.
  Only two numerical suprema were interchanged. Positive square roots give the
  pointwise rate convergence. For $R_n = R_0+n$ and $\epsilon > 0$, the
  increasing open sets $\{s:\lambda_{\Psi,s}-\lambda_{R_n,s}<\epsilon\}$ cover
  compact $I$. A finite subcover lies in the largest one, proving uniform
  convergence. Monotonicity extends it to all real $R \to \infty$. Therefore,
  $\int_I\lambda_{R,s}\ud s\uparrow I_\Psi$.

  The difference of the limiting rate and each masked rate is continuous and
  nonnegative. Once the largest set in a finite subcover contains all of $I$,
  its error is below $\epsilon$ everywhere. Larger radii only decrease that
  error. This is a deterministic compactness argument for rates, not
  convergence of unmasked random flows.

  \medskip
  \noindent\textit{Step 3.~}We take the supremum of fixed-mask moment lower
  bounds.
  At each fixed $R$ use this very same mask in the operator path and mean when
  applying Proposition~\ref{P:raw-block-original-lower}. It gives the lower
  bound with $2\int_I\lambda_{R,s}\ud s-v_\Psi/2$ for the same original point
  liminf. Taking the numerical supremum over these individual bounds proves
  \eqref{E:variational-unmasked-point-lower}. There is no $R(T)$, unmasked noisy
  wave, uniform-in-mask Gaussian estimate or exchange of source and time limits.
  This proves Proposition~\ref{P:variational-remove-mask}.
\end{proof}

\subsection{Full-class dilation and finite variational values for nonnegative controls}

The bound~\eqref{E:variational-unmasked-point-lower} holds for each $\Psi \in
\mathcal{C}_{\mathrm{pos}}$ separately. Before we take the supremum over $\Psi$,
we record how $h_{\Psi,s}$, $v_\Psi$, $\lambda_{\Psi,s}$ and $I_\Psi$ change
under a spatial dilation of the control. The proof of
Proposition~\ref{P:variational-positive-scores} uses this to relate the
variational values for different coefficients $c$.

\LeanMarker{L:variational-control-dilation}
\begin{lemma}[Radon and Rayleigh dilation]\label{L:variational-control-dilation}
  For $r > 0$ the map
  \begin{equation}\label{E:variational-control-dilation}
    (D_r\Psi)(s,x)=r^{d+2-\alpha}\Psi(s,rx)
  \end{equation}
  is a bijection of $\mathcal{C}_{\mathrm{pos}}$, with inverse $D_{1/r}$. With
  the same original covariance,
  \begin{equation}\label{E:variational-dilation-values}
    h_{D_r\Psi,s}(x)=r^2h_{\Psi,s}(rx),\qquad
    v_{D_r\Psi}=r^{4-\alpha}v_\Psi,\qquad
    \lambda_{D_r\Psi,s}=r\lambda_{\Psi,s},\quad
    I_{D_r\Psi}=rI_\Psi.
  \end{equation}
\end{lemma}

\begin{proof}
  We apply the homogeneous Radon pushforward to the covariance mean and the
  Lebesgue Jacobian to its cost. We then dilate the unit Rayleigh tests to
  identify the rate and its time integral.


  \medskip
  \noindent\textit{Step 1.~}We scale the original mean and covariance cost.
  Positive dilation preserves the full smooth positive nonzero class;
  multiplication of its amplitude powers gives the stated inverse. Homogeneity
  of the physical Radon measure gives, for every nonnegative measurable $G$,
  \[
    \int G(ry)\Gamma(\ud y)
    =r^{\alpha-d}\int G(z)\Gamma(\ud z).
  \]
  Indeed, its left pushforward measure on a Borel set $A$ is $\Gamma(A/r) =
  r^{\alpha-d}\Gamma(A)$. Positive Tonelli applies this identity inside the
  temporal convolution, giving $h_{D_r\Psi,s}(x) = r^2h_{\Psi,s}(rx)$. Now use
  the \emph{Lebesgue} substitution $z = rx$ in the original pairing:
  \[
    v_{D_r\Psi}=r^{d+2-\alpha}r^2r^{-d}v_\Psi
    =r^{4-\alpha}v_\Psi.
  \]
  This distinguishes the Radon pushforward from the spatial volume factor; it
  includes singular covariances and white noise.

  The mean has one amplitude factor $r^{d+2-\alpha}$ and one Radon factor
  $r^{\alpha-d}$, whose product is $r^2$. In the cost pairing, multiplying by
  that mean adds $r^2$, while the spatial Lebesgue substitution contributes
  $r^{-d}$. This accounts separately for every exponent in the cost identity.

  \medskip
  \noindent\textit{Step 2.~}We scale the full set of Rayleigh tests.
  The map $U_rf(x) = r^{d/2}f(rx)$ is a unitary bijection on real $L^2$,
  preserving $H^1,H^2$ as sets under its inverse. Weak derivatives give
  $\|\nabla U_rf\|_2^2 = r^2\|\nabla f\|_2^2$ and $\int h_{D_r\Psi,s}(U_rf)^2 =
  r^2\int h_{\Psi,s}f^2$. The real Rayleigh bijection proves the squared rate
  identity; the rates are positive, so their square roots and integrals give
  \eqref{E:variational-dilation-values}. On $H^2$ this is also $H_{D_r\Psi,s}U_r
  = r^2U_rH_{\Psi,s}$. Time has not been rescaled. For compact masks the
  corresponding identity uses $m_R(rx) = m_{R/r}(x)$; a fixed mask radius is not
  invariant. This proves Lemma~\ref{L:variational-control-dilation}.
\end{proof}

\LeanMarker{P:variational-positive-scores}
\begin{proposition}[Intrinsic variational lower bounds]
  \label{P:variational-positive-scores}
  For $c > 0$, the set of variational values
  $\{cI_\Psi-v_\Psi/2:\Psi\in\mathcal{C}_{\mathrm{pos}}\}$ is nonempty and
  bounded above. Its finite real supremum $S_c^{\mathrm{pos}}$ is positive and
  satisfies
  \begin{equation}\label{E:variational-score-homogeneity}
    S_c^{\mathrm{pos}}=c^{(4-\alpha)/(3-\alpha)}S_1^{\mathrm{pos}},
    \qquad
    Q_2\ge S_2^{\mathrm{pos}},
  \end{equation}
  where $Q_2$ is the original point liminf. The coefficient-one value is an
  auxiliary variational value; no first-moment or source-constant identification
  is asserted.
\end{proposition}

\begin{proof}
  We first construct a positive coefficient-one value and scale the entire
  class of values. We use the independently proved point upper bound to show
  these sets are bounded above before taking their finite real suprema.


  \medskip
  \noindent\textit{Step 1.~}We exhibit a positive value.
  Write $k = 3-\alpha > 0$, $q = 4-\alpha = k+1$ and $m = q/k > 1$. Fix one
  $\Psi_0$ from the nonempty class and put $I_0 = I_{\Psi_0} > 0$, $v_0 =
  v_{\Psi_0} > 0$. Positivity of the latter is
  Lemma~\ref{L:compact-raw-profile}. For $r_0 = (I_0/v_0)^{1/k} > 0$, which need
  not be less than one, the exact real-power identity gives
  \[
    I_{D_{r_0}\Psi_0}-v_{D_{r_0}\Psi_0}/2
    =r_0I_0-r_0^qv_0/2=r_0I_0/2>0.
  \]
  The definition of $r_0$ gives $r_0^k=I_0/v_0$. Since $q=k+1$, it follows
  that $r_0^qv_0=r_0I_0$. Thus, subtracting half the cost leaves exactly half
  of the positive linear term.
  This positive value is obtained from a smooth nonnegative control.


  \medskip
  \noindent\textit{Step 2.~}We scale the whole coefficient-one set
  bijectively.
  For $c > 0$ put $r_c = c^{1/k}$ and $K_c = c^m > 0$.
  Lemma~\ref{L:variational-control-dilation} and the exact scalar powers give,
  for every $\Psi$,
  \begin{equation}\label{E:variational-score-set-scaling}
    cI_{D_{r_c}\Psi}-v_{D_{r_c}\Psi}/2
    =K_c(I_\Psi-v_\Psi/2).
  \end{equation}
  The two scalar factors agree: $cr_c=c^{1+1/k}=c^{q/k}=K_c$ and
  $r_c^q=c^{q/k}=K_c$. Consequently, dilation multiplies the linear reward and
  the quadratic cost by the same factor for this choice of $r_c$.
  Since $D_{r_c}$ is a bijection of the entire class, the whole set of
  variational values with coefficient $c$ is $K_c$ times the whole
  coefficient-one set of variational values.


  \medskip
  \noindent\textit{Step 3.~}We bound the sets using the original point upper
  estimate.
  We prove boundedness before taking finite-real suprema. The independent
  Theorem~\ref{T:full-upper}, for the \emph{same canonical solution}, gives
  \[
    \E|u(T,0)|^2\le e(T+2)^d e^{C_{\mathrm{up}}T^\kappa},
    \qquad T\ge1,\quad C_{\mathrm{up}}<\infty.
  \]
  Its prefactor is negligible: choose $\epsilon = \kappa/2 > 0$, use $\log T \le
  (T^\epsilon-1)/\epsilon$ and $\log(T+2) \le \log T+\log3$ for $T \ge 1$, and
  divide by $T^\kappa$. Thus, the original point limsup is at most
  $C_{\mathrm{up}}$. One fixed test/mask lower bound gives a finite lower bound
  for $Q_2$, so $Q_2$ is finite. By~\eqref{E:variational-unmasked-point-lower}
  every coefficient-two variational expression is at most this same $Q_2 \le
  C_{\mathrm{up}}$, uniformly in $\Psi$.
  Equation~\eqref{E:variational-score-set-scaling} at $c = 2$ consequently
  bounds all coefficient-one variational values by $C_{\mathrm{up}}/K_2$. The
  nonempty set has a finite real supremum $S_1^{\mathrm{pos}} > 0$ by the
  positive test above.


  \medskip
  \noindent\textit{Step 4.~}We take finite real suprema and retain their
  interpretation.
  For general $c$, set equality and $K_c > 0$ show that $K_cS_1^{\mathrm{pos}}$
  is the supremum: it is an upper bound, and every smaller number is exceeded by
  $K_c$ times a coefficient-one variational expression within the corresponding
  positive distance of its supremum. This proves finiteness, positivity and the
  first identity in~\eqref{E:variational-score-homogeneity}. Taking the
  numerical supremum of the coefficient-two lower bounds proves the second.
  Although $2^m > 2$, the resulting $S_2^{\mathrm{pos}} > 2S_1^{\mathrm{pos}}$
  is only a comparison of auxiliary variational values. It is not a Jensen
  improvement without a first-moment identification. This proves
  Proposition~\ref{P:variational-positive-scores}.
\end{proof}

For $\Gamma = 0$, the mean, original variance and Rayleigh rate are zero. Set
all these variational values to zero; the free canonical solution $u = 1$ has
the corresponding zero moment rate. Strict positivity above is restricted to
nonzero covariance. Matching an upper variational rate and identifying an exact
pointwise limit remain later steps.
